\subsection{Elements of Riemannian geometry}
\label{sec:elemRiem}
 
We introduce here some elements of Riemannian geometry that will be used throughout this work. Indeed, we will typically assume further in this article that we observe data sampled from a Riemannian manifold, as well as the values of a continuous filter function defined on the points. Although these are standard results, we recall some proof elements in Appendix~\ref{apdx:proofs} for the sake of completeness. A general presentation of the results provided in this section, and notably a proof of Theorem~\ref{thm:bishop_gromov}, can be found in~\cite{petersen2006riemannian}.

We call Riemannian manifold every $C^{\infty}$-manifold $\man$ together with a Riemannian metric $\g$. The Riemannian metric consists of an Euclidean inner product $\g_p$ on each of the tangent spaces $\tangent_p\man$ of $\man$, which satisfies that $p\mapsto\g_p(X_p,Y_p)$ is smooth whenever $X$ and $Y$ are two smooth vector fields defined on $\man$. 

Given some local coordinates $x(p)=(x_1,...,x_d)$ on an open set $U$ of $\man$, 
they induce a basis of the tangent space $\tangent_p\man$. %for all $p\in U$. 
%These are the coordinate vector fields $(\partial_i)_{i=1}^d$.
We denote the associated {\it coordinate vector fields} as $(\partial_i)_{i=1}^d$, and their dual $1$-forms as $(dx_i)_{i=1}^d$.
% so as to have:
%$$dx_i(\partial_j)=
%\begin{cases}
%1, \text{ if } i=j\\
%0, \text{ otherwise}
%\end{cases}$$
%as the vector fields defined by 
%$$\partial_i=\frac{\partial}{\partial x_i}.$$
Under these notations, we can write the Riemannian metric $\g$ as:
$$\g=\sum_{i,j}\g(\partial_i,\partial_j)\cdot dx_i\cdot dx_j,$$
where $dx_i\cdot dx_j$ is the bilinear form: $(v,w)\mapsto dx_i(v)\cdot dx_j(w).$
As such, $\g$ can be represented in local coordinates as a symmetric positive-definite matrix $\G(p)$ with entries parametrized over $U$.

We will denote
$$\lambda(p)=\min\{\sqrt{\rho},\,\rho\in\mathrm{Sp}(\G(p))\},$$
and
$$\mu(p)=\max\{\sqrt{\rho},\,\rho\in\mathrm{Sp}(\G(p))\},$$
where $\mathrm{Sp}(\G(p))$ is the spectrum of $\G(p)$. Notice that since the Riemannian metric is continuous, both $\lambda$ and $\mu$ are continuous as well.

The following proposition states that the Riemannian metric is bi-Lipschitz equivalent to the Euclidean metric in local coordinates.

\begin{proposition}\label{prop:g_equiv}
Let $U$ be an open set in $\man$, on which we are given coordinates $x(p)=(x_1,...,x_d)$. 
%
For every $p\in\man$ and for every $v\in\tangent_p\man$:

$$\lambda(p)\cdot \sqrt{\g_0(v,v)}\leq \sqrt{\g(v,v)}\leq \mu(p)\cdot \sqrt{\g_0(v,v)},$$

where $\g_0=\sum_{i}dx_i\cdot dx_i$ is the Euclidean metric in the local coordinates.
\end{proposition}

Riemannian manifolds can be given a metric space structure as it is possible to measure the ``length" of piecewise smooth curves using the Riemannian metric. This allows to define the {\it geodesic} distance $\dis$ on $\man$. 
The geodesic distance is locally bi-Lipschitz equivalent to the Euclidean distance in local coordinates.

\begin{proposition}\label{prop:dis_equiv}Let $p\in\man$. For every small enough neighborhood $U$ of $p$, we have that for every $q\in U$:
%\begin{equation}
%\label{eq:gamma0}
$$\lambda_0\cdot\dis_0(p,q)\leq \dis(p,q)\leq\mu_0\cdot\dis_0(p,q),$$
%\end{equation}
where $\dis_0(p,q)$ is the Euclidean distance between the representations of $p$ and $q$ in local coordinates and where
$$\lambda_0=\inf_{r\in U}\lambda(r) \textrm{ and }
\mu_0=\sup_{r\in U}\mu(r).$$
In particular $\lambda_0$ tends to $ \lambda(p)$ and
$\mu_0$  tends to  $\mu(p)$ as $\diam(U)$ tends to $0$.
%Moreover, $$\frac{\dis(p,q)}{\dis_0(p,q)}\rightarrow 1\text{ as }q\rightarrow p.$$
% \bert{ca grince un cette convergence de $\lambda_0, \mu_0$, je te propose plutot (a verifier et peut etre à positionner plutot en appendice ?) : Moreover, it is possible to choose $\lambda_0$ and $\mu_0$ continuously in the following sense: for any $r$ small enough  there exists $\lambda_0(r)$ and $\mu_0(r)$ satisfying \eqref{eq:gamma0} for any $q \in U \cap B_0(p,r)$, such that $\lambda_0(r)$ and $\mu_0(r)$ tends to $\lambda_0$ and $\mu_0$ as $r$ tends to zero.}
\end{proposition}











%The above result shows in particular that the manifold topology is equivalent to the metric topology induced by the geodesic distance.\\
The Riemannian metric also induces a Borel measure $\vol$ called the {\it volume measure}. It can be first defined by specifying the expected value of a function $f\colon\man\rightarrow\mathbb{C}$ compactly supported on a single chart $\varphi\colon U\subseteq\man\rightarrow V\subseteq\mathbb{R}^d$:

$$\int_{\man}f\,d\vol=\int_{V}\left(f\cdot\sqrt{\deter(\G)}\right)\circ\varphi^{-1}\,d\lambda,$$

where \begin{align*}
\G\colon U&\longrightarrow\mathbb{R}^{d\times d}\\
p&\longmapsto(\g(\partial_{i_{|p}},\partial_{j_{|p}}))_{i,j}
\end{align*}
and $\lambda$ is the Lebesgue measure on $\mathbb{R}^d$.\\
It can be checked, with the help of the substitution rule, that the above definition does not depend on the choice of the coordinate neighborhood.\\
This definition can be extended to general compactly supported functions by using a partition of unity of a coordinate neighborhood cover of $\man$: we simply sum the expected values over the elements of the partition.\\
Subsequently, as the expected value of compactly supported functions is well defined, the Riesz representation theorem allows to consider the unique associated positive Borel measure on $\man$: this is the volume measure $\vol$.\\\\
 We have the following asymptotic equivalence result for the volume of small geodesic balls.
\begin{proposition}{Volume of small geodesic balls.}\label{prop:geo_vol}\\
Let $\ball(p,\eps)$ be the geodesic ball around $p\in\man$ of radius $\eps$. We have that:
$$\vol\left(\ball(p,\eps)\right)\underset{\eps\rightarrow 0}{\sim} \alpha_d\cdot\eps^d,$$
where $\alpha_d$ is the volume of the unit ball in $\R^d$.
%In particular, there exists $\epsilon'>0$ and $\lambda_1,\mu_1>0$ such that for every $\epsilon\leq\epsilon'$:
%
%$$\lambda_1\cdot\epsilon^d\leq \vol\left(\ball(p,\epsilon)\right)\leq\mu_1\cdot\epsilon^d.$$
\end{proposition}

%For a power series expansion of the volume of small geodesic balls to higher orders, see~\cite{gray1974volume} under the assumption that the manifold is analytic, which is more restrictive than our case.\bert{on retire ce commentaire ?}\\
For the following result we will assume a lower bound on the {\it Ricci curvature} of $\man$. The Ricci curvature $\ricci$ is a symmetric bilinear form on the tangent spaces and can be thought of as the Laplacian of the metric $\g$, see for instance Chapter 9 in~\cite{petersen2006riemannian}. We will adopt the convention that $\ricci\geq k$ means that all eigenvalues $\rho$ of $\ricci$ satisfy $\rho\geq k$. 

The next theorem allows to compare the volume of geodesic balls in $\man$ to the volume of geodesic balls in certain model manifolds $\spaceform^d_k$ called {\it constant-curvature space forms}, $d$ being the dimension of both $\man$ and $\spaceform^d_k$ (see for example Subsection 9.1. in~\cite{petersen2006riemannian}).
Denote $v(d,k,r)$ as the volume of a ball of radius $r$ in $\spaceform^d_k$. 

\begin{theorem}{Bishop-Cheeger-Gromov Theorem, see for instance Lemma~36 in~\cite{petersen2006riemannian}.}\\
\label{thm:bishop_gromov}
Suppose that $\man$ is a $d$-dimensional complete Riemannian manifold such that $\ricci\geq(d-1)\cdot k$, for some $k \in \R$. Then  for any $p \in \man$, 
$$\eps\longmapsto\frac{\vol\left(\ball(p,\eps)\right)}{v(d,k,\eps)} $$
is a nonincreasing function whose limit is 1 as $\eps\rightarrow 0$. 
\end{theorem}
 
Note that in the above theorem, $v(d,k,\eps)$ is independent of the base point $p$. %of $\ball(p,r)$. 
This is very convenient for providing bounds on the volume of geodesic balls in $\man$ that are uniform with respect to the choice of base points. Notice also that since $\spaceform^d_k$ is a $d-$dimensional Riemannian manifold, Proposition~\ref{prop:geo_vol} applies for $v(d,k,\eps)$, i.e.,
$$v(d,k,\eps)\underset{\eps\rightarrow 0}{\sim} \alpha_d\cdot \eps^d.$$

We now define the gradient of a smooth real valued function on a Riemannian manifold. 
%Like mentioned above, 
Let $x(p)=(x_1,...,x_d)$ be some local coordinates on an open set $U$ of $\man$
%induce a basis of the tangent spaces $\tangent_p\man$ for all $p\in U$. These are the 
with associated coordinate vector fields $(\partial_i)_{i=1}^d$.
Let $f:\man\rightarrow\R$ be a smooth function. We can define its differential $df$ as the $1$-form that satisfies:
$df(\partial_i)=\frac{\partial f}{\partial x_i}$, 
%\mathieu{Shouldn't it be $\frac{\partial \hat f}{\partial x_i}$ where $\hat f = f\circ \varphi^{-1   }$ is the representative of $f$ in the local coodinates?}\ziy{Oui, la notation $\frac{\partial f}{\partial x_i}$ veut dire ce que tu as écrit. Je crois que vu qu'on note $x_i$ le choix de coordonnées n'est pas ambigu.} 
for all $1\leq i \leq d.$
The gradient $\nabla f$ of $f$ is then defined as the vector field satisfying:
$\g(\nabla f,v)=df(v),$
for all vector fields $v$. It can also be seen as the Riesz representative of $df$.
Given our local coordinates, we can express the gradient vectors in their associated coordinate vector field bases.
\begin{proposition}\label{prop:grad}
In the local coordinates associated to the open set $U$, we have that:
$$\nabla f=\sum_{i=1}^d a_i\cdot\partial_i,$$
where:
$$(a_1,...,a_d)=\left(\frac{\partial f}{\partial x_1},...,\frac{\partial f}{\partial x_d}\right)\cdot\left[\g(\partial_i,\partial_j)\right]_{i,j}^{-1}.$$
In particular, we have:
$$\frac{1}{\mu(p)}\cdot \sqrt{\sum_{i=1}^d\frac{\partial f(p)}{\partial x_i}^2}\leq \sqrt{\g\left(\nabla f(p),\nabla f(p)\right)}\leq \frac{1}{\lambda(p)}\cdot \sqrt{\sum_{i=1}^d\frac{\partial f(p)}{\partial x_i}^2}$$

\end{proposition}












